\par\clearpage
\noindent\null\hfill{\small\textbf{ON THE HANOVER TRIANGULATION.}}\hfill\makebox[0pt][r]{\fontsize{13}{16}\selectfont 411}\par\medskip
\noindent also became apparent that Falkenberg could be connected to Wilsede by
cutting a sightline through the Becklinger Holz, and there was hope that
the two points Hauselberg and Garssen could be connected by way of one or a few
intermediate points and by cutting several sightlines through the woods that
separate them. But as regards such cuts, without which it is altogether
impossible to get through in these parts, I thought I must set myself two rules:
first, to avoid all the finer kinds of timber (such as oak woods) as far as
at all possible, and second, to determine the direction of the cuts
\textls[120]{in advance} with the \textls[120]{utmost precision}, so that they would be as
narrow as possible and there would be no need to fell a single tree unnecessarily.
To achieve this, however, very precise measurements had to be made first, some
of which had to be ingeniously arranged and combined. These could be carried
out accurately and with the least possible expenditure of time only once the
principal measurements themselves had advanced into this area, and I may say
here in advance that I subsequently achieved these aims to my fullest
satisfaction.

I used the time between my return from the reconnaissance trip and the actual
start of the surveying operations to get the instruments fully ready for use,
make all the necessary preparations, and give Captain \textsc{Müller} and
Lieutenant \textsc{Hartmann} the instructions needed to coordinate the operations
perfectly. On 16~June I set out for Lichtenberg, the first triangulation station
to be taken up this year, where my son had already gone ahead with some of the
instruments.

The actual surveying work lasted four months this year, namely until the middle
of October, and during that time work was completed at \textls[120]{nine}
principal triangulation stations: in 1821 only five could be taken up in a period not
much shorter. This much more rapid progress is chiefly due to the increased
number of assistants and heliotropes, but still more important is the resulting
gain in the accuracy of the measurements themselves.\quad\mbox{.\ .\ .\ .\ .\ .}

The purpose of the triangulation, as part of the arc measurement, is to connect
the Göttingen observatory through a continuous system of triangles with the \mbox{Dan-}
\par\smallskip\noindent\hfill\makebox[0.18\linewidth][l]{52\textsuperscript{*}}\par
\clearpage
\noindent\makebox[0pt][l]{\fontsize{13}{16}\selectfont 412}\hfill{\small\textbf{POSTHUMOUS PAPERS.}}\hfill\null\par\medskip
\noindent ish triangles. To this end it was most advantageous to make the triangles
as large as possible and generally to choose triangulation stations at the
highest sites, which afford the widest view. For the most part these points in
themselves are of no immediate interest for the geography of the kingdom. In
addition to my main objective, however, I have everywhere sought to make my
operations as useful as possible for that geography. I carefully determined the
positions of all places visible from more than one principal triangulation
station, some with an accuracy equal to that of the positions of the principal
stations. Among them I name the towns of Hannover, Braunschweig, Celle, Lüneburg,
Neustadt am Rübenberge, and Burgdorf. The number of villages whose positions I
determined is very large. The way has been opened to extend this harvest over a
larger part of the kingdom, or over the whole of it.\quad\mbox{.\ .\ .\ .\ .\ .}

\bigskip
\begin{center}
[From the report of 16~February 1825 ȟber die im Jahre 1824\par
ausgeführten trigonometrischen Arbeiten.«]
\end{center}

\par\noindent\hspace*{\parindent}\mbox{.\ .\ .\ .\ .\ .}\quad Results: The report and the accompanying map show that the work
of 1824 offers several connections from the triangles of the earlier years as
far as Bremen. The simplest [comprising the triangles Wilsede-Falkenberg-Elmhorst,
Wilsede-Elmhorst-Litberg, Wilsede-Litberg-Hamburg, Wilsede-Litberg-Zeven,
Wilsede-Zeven-Steinberg, and Zeven-Steinberg-Bremen\,*)] is drawn with heavy
solid lines, and two new triangles [Zeven-Bremen-Brillit and Brillit-Bremen-Garlste]
are also attached to the side Bremen-Zeven. Had it been possible to find that
simplest system at the outset, the other triangles drawn with fine solid lines
[with Bullerberg, Bottel, and Brüttendorf as vertices] could indeed have been
dispensed with entirely. But the report shows how many difficulties had to be
overcome before the plan for that system could be worked out, and it would
likewise have been quite impracticable to cut the various necessary sightlines
precisely and quickly without first having acquired
\par\vfill\noindent\rule{0.25\linewidth}{0.4pt}\par\smallskip
{\footnotesize\noindent\hspace*{\parindent}[*) See the sketch of the triangles on p.~299].\par}
\clearpage
\noindent\null\hfill{\small\textbf{ON THE HANOVER TRIANGULATION.}}\hfill\makebox[0pt][r]{\fontsize{13}{16}\selectfont 413}\par\medskip
\noindent a very accurate knowledge of the positions of the sites.
In the present mathematical state of higher geodesy, moreover, even the latter
triangles, those drawn with fine lines, are by no means to be considered
superfluous. Indeed, taking them into account according to quite definite
principles must help increase the precision of the final results.\quad\mbox{.\ .\ .\ .\ .\ .}

Finally, it is also of great importance that the three triangles between the
five points Falkenberg, Elmhorst, Wilsede, Litberg, Hamburg have provided a new
connection from the southern triangles in the Kingdom of Hanover as far as
Hamburg. It is preferable to the earlier, far more complicated one [via
Hauselberg, Wulfsode, Timpenberg, Nindorf, and Lüneburg] and will therefore,
according to the principles indicated above, double the accuracy of the
results.\quad\mbox{.\ .\ .\ .\ .\ .}

\bigskip
\begin{center}
[From the report of 21~November 1827, »betreffend die weitere Ausdehnung\par
der Gradmessungsarbeiten.«]
\end{center}

\par\noindent\hspace*{\parindent}\mbox{.\ .\ .\ .\ .\ .}\quad In 1821---1823, 11000 thalers were spent on measuring the chain of
triangles as far as Hamburg, and in 1824, 1825 about 7000 thalers on the chain
of triangles extending westward from there as far as Ostfriesland. From these
costs, however, one must deduct the expenditure on acquiring instruments and
bringing the English zenith sector from Altona to Göttingen, which may have
amounted to between 2500 and 3000 thalers, so that the costs of the triangulation
proper may have been about 15000 thalers. Now, according to the overview map,
the parts of the country still untouched, taken together, seem to be not much
larger than the area already covered by triangles. For all my uncertainty about
the difficulties of the terrain, it is scarcely likely that they could be
greater than those I had to contend with especially in 1822 and 1824. If I
further point out that the operations whose main purpose is to improve the
geography of the country, even when carried out to a standard worthy of the
undertaking, do not require \textls[120]{the} degree of extreme precision in
measurement demanded by an arc measurement proper, then it seems
\par\clearpage
\noindent\makebox[0pt][l]{\fontsize{13}{16}\selectfont 414}\hfill{\small\textbf{POSTHUMOUS PAPERS.}}\hfill\null\par\medskip
\noindent that there are grounds for hoping the extension of the triangulation
over the parts of the kingdom still untouched can be paid for with a smaller
sum, perhaps with 12000 thalers\,*).

To make such a triangulation as useful as possible for improving the geography,
one must go beyond laying out networks of principal triangles of the first
order over the parts of the country in question and combine this work with
determining the positions of as many secondary points as possible, especially
those that allow precise determinations and usually last for centuries, hence
especially church towers. I made this consideration a duty even in the earlier
surveys, although it had to remain subordinate to the main objective, and the
number of points determined in those surveys already exceeds 500.\quad\mbox{.\ .\ .\ .\ .\ .}

This statement in numbers of the positions of a large number of fixed points
(namely how far north or south, west or east of any chosen starting point,
e.g.\ the Göttingen observatory), \textls[120]{accurate to within a few feet},
must be regarded as the principal topographical result of the operations. It
retains its value for centuries, since the majority of the points remain even
if some are lost over time, and any changes that may arise as a result can
readily be made good. It provides a sound basis for all detailed surveys: all
the uncertainties that make surveys without such fixed reference points more
difficult, distort them, and make it impossible to combine them into an error-free
whole, disappear completely. This puts careless workers under a strict and
unavoidable check. Each plane-table sheet is worked on independently of the
others, so no error propagates to other sheets. Finally, all the individual
sheets fit together of their own accord to form a precisely oriented and
everywhere consistent whole. It is clear that the great costs invariably
entailed by detailed surveys of considerable extent must be greatly reduced
by such a reliable procedure, but this reliable procedure is not the only thing
\par\vfill\noindent\rule{0.25\linewidth}{0.4pt}\par\smallskip
{\footnotesize\noindent\hspace*{\parindent}*) The costs required in the current year 1827 do not enter into this reckoning, as they relate to the purely astronomical part of the arc measurement.\par}
\clearpage
\noindent\null\hfill{\small\textbf{ON THE HANOVER TRIANGULATION.}}\hfill\makebox[0pt][r]{\fontsize{13}{16}\selectfont 415}\par\medskip
\noindent that speeds up the work. It is also very important in this regard
that the use of the measuring chain thereby becomes almost entirely superfluous
and is needed only in exceptional cases, since the triangulation already
provides the baselines of itself, and with a precision that the ordinary chain
could not achieve at all.

But even where detailed surveys already exist, as in most of the administrative
districts of the kingdom as it formerly stood, the fixed points provide the
means to correct the errors that arose in putting them together. In this way,
even maps based on imperfect surveying methods, if otherwise good in their
fine detail, can be reworked into representations that meet higher demands.

As for the measures to make the triangulations as useful as possible for
improving the geography, a distinction must accordingly be made between the
surveys already carried out and those that may in future be extended over
other parts of the country.

For the latter, obtaining precise positions for as many fixed points as
possible will have to be treated as a main objective from the outset.

In the surveys already carried out, on the other hand, this consideration
was indeed regarded as subordinate. Yet since, as already mentioned, I always
kept it in view to the extent possible without impeding the main undertaking,
the northern and southern regions must again be distinguished here as regards
the results.

In the northern (larger) half, i.e.\ roughly from the town of Hildesheim to
the sea, and hence in the flat part of the country, the yield in the respect
mentioned has been so abundant that little or nothing is left to be desired.\quad\mbox{.\ .\ .\ .\ .\ .}

In the southernmost part of the kingdom, on the other hand, the number of
church towers whose positions have been precisely determined is much smaller,
since, partly because of the size of the triangles and partly because of the
mountainous character of the country, only a few towers could be seen from
more than one principal triangulation station at the same time. To improve the geography of
the kingdom, and particularly to provide a similarly sound basis for a detailed
survey of the southern parts of the Hildesheim territory and the Eichsfeld,
it would therefore \mbox{cer-}
\par
